%-----------------------------------------------------------------------------%
\chapter*{Extended Abstract}
\addcontentsline{toc}{chapter}{Extended Abstract}
%-----------------------------------------------------------------------------%

\renewcommand{\thetable}{\arabic{table}}
\renewcommand{\thefigure}{\arabic{figure}}
\setcounter{figure}{0}
\setcounter{table}{0}

\vspace*{0.2cm}
{
\setlength{\parindent}{0pt}

\begin{center}
{\large\textbf{EXTENDED ABSTRACT}}

\medskip
{\normalsize\textbf{Undergraduate Thesis}}

\medskip
{\normalsize\textbf{\textit{\judulInggris}}}

\medskip
\textbf{Student Name} : \penulis \\
\textbf{Student ID Number (NPM)} : 2206051462 \\
\textbf{Study Program} : \program
\end{center}

\bigskip

% #20 opening gap sentence added
Every day, review platforms and social media generate enormous volumes of
user-generated text. Automatically extracting sentiment and thematic content
from this data has become a central challenge in Natural Language Processing
(NLP). Conventional document-level sentiment analysis reduces each document
to a single polarity label, which overlooks the fact that a review often
discusses multiple aspects with differing opinions. Yet existing joint
sentiment-topic approaches either require aspect-level supervision or rely
on pipeline architectures that do not allow the two objectives to reinforce
each other during training. This limitation motivates end-to-end joint
approaches that simultaneously model sentiment and topics, offering richer
output without requiring aspect-level annotations.

% #17 JDST-DEC name bolded on first mention; #3 "summarised"; #4 loss weight extended; #5 K notation
To address this, the present study proposes a \textit{Joint Document-Level
Sentiment-Topic} model based on \textit{Deep Embedded Clustering}
(\textbf{JDST-DEC}) within a \textit{Multi-Task Learning} (MTL) framework.
The model operates in two stages. In the first stage, the 768-dimensional
\texttt{[CLS]} representation extracted by a frozen BERT encoder, which
reduces computational cost and preserves stable pre-trained
representations, is fed into a symmetric fully connected autoencoder
with architecture $768 \to 2022 \to 2022 \to 256$, trained using mean
squared error (MSE) to learn a compact 256-dimensional latent
representation. In the second stage, this latent space is jointly optimised
for sentiment classification via cross-entropy loss
($\mathcal{L}_{\text{sent}}$) and topic clustering via DEC with KL
divergence loss ($\mathcal{L}_{\text{clust}}$), under the joint objective:
\[
    \mathcal{L} = \eta\,\mathcal{L}_{\text{sent}} +
    \gamma\,\mathcal{L}_{\text{clust}}, \qquad \eta = \gamma = 0{.}5,
\]
where equal weighting balances the sentiment classification and clustering
objectives without favouring either. Cluster centroids are initialised
using \textit{k}-means on the pre-trained latent space. The overall
research workflow is summarised in Figure~1.

% #2 caption no longer has "Figure 1." prefix
\begin{figure}[H]
    \centering
    \includegraphics[width=0.82\textwidth]{assets/pics/other/Flowchart.png}
    \caption{Schematic diagram of the research framework.}
    \label{fig:ea_workflow}
\end{figure}

% #1 fairness statement; #5 K notation; #18 MTL-GUI consistent
Experiments were conducted on two benchmark datasets: IMDB Movie Reviews
\parencite{maas2011} (50{,}000 reviews) and the Yelp Review Dataset
\parencite{zhang2015} (30{,}000 reviews), both with binary sentiment labels.
Three topic cardinalities $K\in\{30,50,80\}$ were evaluated.
\textbf{JDST-DEC} was compared against NN-BERT as a single-task sentiment
baseline, MTL-GUI \parencite{gui2022mtml}, BERTopic-HDBSCAN, and
BERTopic-KMeans. All baselines were re-evaluated under identical dataset
splits and evaluation protocols to ensure a fair comparison. Sentiment
performance was measured using accuracy and F1-score; topic quality was
assessed using \textit{Topic Coherence} based on NPMI (TC-NPMI) and
\textit{Topic Diversity} (TD).

Pre-training results show that reconstruction loss drops sharply in the
early epochs and stabilises after epoch~20, reaching final training losses
of approximately $1{.}41 \times 10^{-4}$ on IMDB and $1{.}55 \times 10^{-4}$
on Yelp. Validation loss remains close to training loss throughout, indicating
no overfitting. Selected pre-training loss values are shown in Table~1.

\begin{table}[H]
    \centering
    \caption{Selected pre-training \textit{loss} values on the IMDB and
    Yelp datasets.}
    \label{tab:ea_pretrain}
    \renewcommand{\arraystretch}{1.15}
    \begin{tabular}{ccccc}
        \toprule
        \multirow{2}{*}{\textbf{Epoch}}
          & \multicolumn{2}{c}{\textbf{IMDB}}
          & \multicolumn{2}{c}{\textbf{Yelp}} \\
        \cmidrule(lr){2-3}\cmidrule(lr){4-5}
          & \textit{Train} & \textit{Val.}
          & \textit{Train} & \textit{Val.} \\
        \midrule
         1  & $2{.}370 \times 10^{-3}$ & $2{.}544 \times 10^{-3}$
            & $3{.}790 \times 10^{-3}$ & $4{.}062 \times 10^{-3}$ \\
        10  & $3{.}780 \times 10^{-4}$ & $4{.}120 \times 10^{-4}$
            & $2{.}780 \times 10^{-4}$ & $3{.}040 \times 10^{-4}$ \\
        20  & $1{.}370 \times 10^{-4}$ & $1{.}550 \times 10^{-4}$
            & $1{.}670 \times 10^{-4}$ & $1{.}860 \times 10^{-4}$ \\
        30  & $1{.}410 \times 10^{-4}$ & $1{.}590 \times 10^{-4}$
            & $1{.}550 \times 10^{-4}$ & $1{.}730 \times 10^{-4}$ \\
        \bottomrule
    \end{tabular}
\end{table}

% #6 "percentage points"; #19 consistent number format
For sentiment classification, \textbf{JDST-DEC} consistently outperforms
NN-BERT and MTL-GUI across all values of $K$ on both datasets (Table~2).
On IMDB, \textbf{JDST-DEC} achieves its best F1-score of 0.8986 at $K = 50$,
compared with 0.8612 for NN-BERT and 0.8317 for MTL-GUI, corresponding to
improvements of approximately 3.7 and 6.7 percentage points respectively.
On Yelp, \textbf{JDST-DEC} reaches an F1-score of 0.9020 at $K = 50$,
compared with 0.8715 and 0.8401, representing gains of approximately 3.1 and
6.2 percentage points. Performance remains stable across all values of $K$,
indicating that sentiment classification quality is relatively insensitive
to topic cardinality. Among the compared methods,
\textbf{JDST-DEC} achieves the highest sentiment scores on Yelp and
competitive results on IMDB.

% #16 dagger notation added to method name; #15 identical evaluation settings in footnote
% #16 dagger notation added to method name; #15 identical evaluation settings in footnote
\begin{table}[H]
\centering
\caption{Comparison of sentiment classification results on IMDB and Yelp.}
\label{tab:ea_sentiment}
\small
\renewcommand{\arraystretch}{1.55}
\setlength{\tabcolsep}{6pt}
\begin{adjustbox}{width=\textwidth}
\begin{tabular}{|l|cc|cc|cc|cc|cc|cc|}
\hline
\multirow{2}{*}{Method}
  & \multicolumn{6}{c|}{IMDB}
  & \multicolumn{6}{c|}{Yelp} \\
\cline{2-13}
  & \multicolumn{2}{c|}{}
  & \multicolumn{2}{c|}{Acc \qquad F1}
  & \multicolumn{2}{c|}{}
  & \multicolumn{2}{c|}{}
  & \multicolumn{2}{c|}{Acc \qquad F1}
  & \multicolumn{2}{c|}{} \\
\hline
NN-BERT
  & \multicolumn{6}{c|}{0.8581 \qquad 0.8612}
  & \multicolumn{6}{c|}{0.8664 \qquad 0.8715} \\
\hline
  & \multicolumn{2}{c|}{$K=30$}
  & \multicolumn{2}{c|}{$K=50$}
  & \multicolumn{2}{c|}{$K=80$}
  & \multicolumn{2}{c|}{$K=30$}
  & \multicolumn{2}{c|}{$K=50$}
  & \multicolumn{2}{c|}{$K=80$} \\
\cline{2-13}
  & Acc & F1
  & Acc & F1
  & Acc & F1
  & Acc & F1
  & Acc & F1
  & Acc & F1 \\
\hline
MTL-GUI
  & 0.8327 & 0.8326
  & 0.8317 & 0.8317
  & 0.8334 & 0.8334
  & 0.8385 & 0.8385
  & 0.8380 & 0.8380
  & 0.8402 & 0.8401 \\
\hline
\textbf{JDST-DEC}
  & \cellcolor{green!25}\textbf{0.8957} & \cellcolor{green!25}\textbf{0.8963}
  & \cellcolor{green!25}\textbf{0.8967} & \cellcolor{green!25}\textbf{0.8986}
  & \cellcolor{green!25}\textbf{0.8941} & \cellcolor{green!25}\textbf{0.8930}
  & \cellcolor{green!25}\textbf{0.8965} & \cellcolor{green!25}\textbf{0.8944}
  & \cellcolor{green!25}\textbf{0.8983} & \cellcolor{green!25}\textbf{0.9020}
  & \cellcolor{green!25}\textbf{0.8935} & \cellcolor{green!25}\textbf{0.8953} \\
\hline
\end{tabular}
\end{adjustbox}
\vspace{0.3em}
\noindent{\footnotesize
  \colorbox{green!25}{\,\,} highest value in column;
  
  All baselines evaluated under identical dataset splits and protocols.}
\end{table}

% #7 "which is expected" removed
For topic modelling, BERTopic achieves the highest TD values across most
configurations (Table~3), with BERTopic-KMeans reaching 0.9567 on IMDB
at $K = 30$. For TC-NPMI, BERTopic-HDBSCAN leads at $K = 30$ and $K = 50$
on IMDB, while BERTopic-KMeans leads at $K = 80$. On Yelp, BERTopic
variants achieve the highest TC-NPMI values across all settings.
\textbf{JDST-DEC} produces lower TC-NPMI and TD values than BERTopic, as
joint optimisation prioritises sentiment-discriminative representations
over purely coherent topic structures. Among the joint models,
\textbf{JDST-DEC} consistently outperforms MTL-GUI on both metrics and
both datasets across every value of $K$.

% #16 dagger in topic table; #15 footnote; #17 JDST-DEC bolded
\begin{table}[H]
\centering
\caption{Comparison of topic modelling results on IMDB and Yelp.}
\label{tab:ea_topic}
\small
\renewcommand{\arraystretch}{1.55}
\setlength{\tabcolsep}{6pt}
\begin{adjustbox}{width=\textwidth}
\begin{tabular}{|l|cc|cc|cc|cc|cc|cc|}
\hline
\multirow{3}{*}{Method}
  & \multicolumn{6}{c|}{IMDB}
  & \multicolumn{6}{c|}{Yelp} \\
\cline{2-13}
  & \multicolumn{2}{c|}{$K=30$}
  & \multicolumn{2}{c|}{$K=50$}
  & \multicolumn{2}{c|}{$K=80$}
  & \multicolumn{2}{c|}{$K=30$}
  & \multicolumn{2}{c|}{$K=50$}
  & \multicolumn{2}{c|}{$K=80$} \\
\cline{2-13}
  & TC-NPMI & TD
  & TC-NPMI & TD
  & TC-NPMI & TD
  & TC-NPMI & TD
  & TC-NPMI & TD
  & TC-NPMI & TD \\
\hline
BERTopic-HDBSCAN
  & \cellcolor{green!25}\textbf{0.1714} & 0.7966
  & \cellcolor{green!25}\textbf{0.1823} & 0.7680
  & 0.1751 & 0.7625
  & 0.1092 & 0.7233
  & \cellcolor{green!25}\textbf{0.1666} & \cellcolor{green!25}\textbf{0.7041}
  & \cellcolor{green!25}\textbf{0.1646} & \cellcolor{green!25}\textbf{0.6671} \\
\hline
BERTopic-KMeans
  & 0.1566 & \cellcolor{green!25}\textbf{0.9567}
  & 0.1596 & \cellcolor{green!25}\textbf{0.9000}
  & \cellcolor{green!25}\textbf{0.1874} & \cellcolor{green!25}\textbf{0.8413}
  & 0.1077 & \cellcolor{green!25}\textbf{0.7800}
  & 0.1120 & 0.6740
  & 0.1465 & 0.6438 \\
\hline
MTL-GUI
  & 0.0200 & 0.5833
  & 0.0108 & 0.5800
  & 0.0060 & 0.4913
  & \cellcolor{green!25}\textbf{0.1578} & 0.6600
  & 0.1086 & 0.6720
  & 0.1123 & 0.5000 \\
\hline
\textbf{JDST-DEC}
  & 0.0803 & 0.2256
  & 0.0551 & 0.2100
  & 0.0675 & 0.2875
  & 0.0788 & 0.3267
  & 0.0369 & 0.2280
  & 0.0563 & 0.3150 \\
\hline
\end{tabular}
\end{adjustbox}
\vspace{0.3em}
\noindent{\footnotesize
  \colorbox{green!25}{\,\,} highest value in column;
  
  All baselines evaluated under identical dataset splits and protocols.}
\end{table}

% #8 training dynamics interpretation added
Analysis of the joint training dynamics shows that
$\mathcal{L}_{\text{sent}}$ decreases consistently from early iterations,
while $\mathcal{L}_{\text{clust}}$ tends to increase throughout training.
This reflects the DEC mechanism, in which the target distribution is
iteratively updated based on current cluster assignments, keeping the
KL divergence active and allowing the cluster structure to continue
evolving during joint optimisation. The opposing trends of the two losses
suggest that the model progressively trades raw reconstruction accuracy
for a latent geometry that is simultaneously sentiment-discriminative and
cluster-coherent.

% #9 fewer token examples; #10 "lexical polarity"
Interpretability analysis using occlusion-based token attribution on the
Yelp dataset ($K = 30$) shows that tokens such as \textit{``super''} and
\textit{``average''} carry attributions consistent with their expected
polarity, while certain tokens exhibit signs inconsistent with their
lexical polarity, indicating that the model captures context-sensitive
interactions rather than relying on surface-form cues alone.

% #11 frozen encoder limitation added
\paragraph{Limitations.}
The current study has several limitations. The experiments are limited to
binary sentiment classification on English-language datasets; performance
on multi-class or multilingual settings remains to be established. The
use of a frozen BERT encoder, while computationally efficient, may limit
the model's ability to adapt its representations to domain-specific
sentiment signals. The autoencoder architecture and DEC objective are
fixed throughout training, and it is unclear whether deeper or alternative
architectures would improve topic quality without sacrificing sentiment
accuracy. Finally, evaluation relies on automated topic metrics
(TC-NPMI, TD), which do not fully capture human judgements of topic
usefulness.

% #12 encoder fine-tuning added to future work
In future work, these limitations motivate exploring encoder fine-tuning
strategies that balance adaptation with stability, deeper autoencoder
architectures and alternative clustering objectives to improve topic
diversity, and extensions to multilingual and multi-class sentiment
settings. Incorporating aspect-level evaluation would further
characterise the granularity of the learned topic representations.

% #13 impact sentence added
In conclusion, \textbf{JDST-DEC} demonstrates an effective mechanism for
joint sentiment-topic learning: integrating DEC into an MTL framework
with a frozen BERT encoder improves sentiment classification over both
single-task and existing joint-model baselines, and produces interpretable
topic clusters, all without requiring aspect-level annotations. By showing
that a single latent space can serve both objectives simultaneously, this
work opens a practical path toward richer, annotation-efficient document
understanding in large-scale review and social-media settings.

\bigskip
% #14 Topic Modeling keyword added
\noindent\textbf{Keywords}: \textit{Deep Embedded Clustering},
\textit{Joint Document-Level Sentiment-Topic Modeling},
\textit{Multi-Task Learning}, \textit{Topic Modeling},
\textit{BERT}.

}

% Restore chapter-prefixed numbering after abstract
\renewcommand{\thetable}{\thechapter.\arabic{table}}
\renewcommand{\thefigure}{\thechapter.\arabic{figure}}

\newpage